\section{Conclusion}
\label{sec:conclusion}

HMRS addresses how an agent finds applicable evidence in a changing,
multi-domain document corpus. It makes the hierarchy used to manage the
corpus available at query time: content-informed placement and descriptions
help the agent choose a collection and document, and source references let
it inspect the returned text. The hierarchy makes scope an explicit choice,
not a replacement for lexical or vector retrieval and not a guarantee that
every call stays within the intended scope.

The evaluation shows where this design helps and what it does not yet
establish. A rename-and-move probe exposed stale index entries, so correct
filing did not guarantee current index visibility. In the routing tests,
the configuration with descriptive cards reached the labelled collection
more often than the nominally name-only configuration, although those
runs do not isolate the effect of descriptions. Correct routing still
left some evidence packs without the literal target string. Combining
scored document text with an extra vector search and prioritised windows
improved containment relative to the reference configuration across three
runs on the same questions. In the fixed-pool replay, ordering by verifier
score instead of fused rank reduced top-five gold-document hits. Neither
literal containment nor a top-five hit establishes support for an answer.

The blind deployment comparison found no statistically significant paired
difference after correction. The deployments differed in search scope and
source access, so this was not a matched test of answer quality. The supported conclusion is
narrower: HMRS gives operators and agents separate ways to check document
visibility, where the agent searched and what text it reached. A next test
should measure claim-level support on independent questions under matched
search scopes and reading budgets.
